%================================================
% PACKAGES AND THEMES
%================================================
\documentclass[t,aspectratio=169,xcolor=dvipsnames]{beamer}

\usetheme{SimplePlusAIC}

\usepackage{graphicx}
\usepackage{booktabs}
\usepackage{tcolorbox}
\usepackage{tikz}
\usetikzlibrary{arrows.meta, shapes.geometric, shapes.misc, fit, positioning, calc}
\usepackage{makecell}
\usepackage{setspace}
\usepackage[T1]{fontenc}
\usepackage{helvet}
\usepackage{amsmath}
\usepackage{bm}
\usepackage{ragged2e}
\usepackage[svgnames,table]{xcolor}
\arrayrulecolor{black}
\setlength{\arrayrulewidth}{0.20mm}
\renewcommand{\arraystretch}{1.35}

\newcommand{\tem}[1]{\textbf{\textcolor{red}{#1}}}

%================================================
% TITLE PAGE
%================================================
\title[Workshop Sesi DL]{Pemrosesan Bahasa Alami}
\subtitle{Deep Learning for NLP — From RNN to Transformers}
\author[SD25-32202]{Martin Manullang}
\institute[ITERA]{Institut Teknologi Sumatera \\ mctm.web.id}
\date{\textcolor{nyublue}{2026}}

%================================================
% AT BEGIN SECTION - SECTION SEPARATOR
%================================================
\AtBeginSection[]{
  \begin{frame}[plain,c]
    \begin{center}
      \begin{tcolorbox}[
        colback=nyublue!10!white,
        colframe=nyublue,
        boxrule=1.5mm,
        arc=3mm,
        halign=center,
        valign=center,
        height=0.45\textheight,
        width=0.85\textwidth
      ]
        \Huge\textbf{\textcolor{nyublue}{\insertsectionhead}}
      \end{tcolorbox}
    \end{center}
  \end{frame}
}

%================================================
% BEGIN DOCUMENT
%================================================
\begin{document}

\begin{frame}
    \titlepage
\end{frame}

\begin{frame}
    \frametitle{OUTLINE}
    \framesubtitle{Peta materi pertemuan ini}
    \small
    \begin{spacing}{1.15}
        \tableofcontents
    \end{spacing}
\end{frame}

%================================================
% SLIDE PENDAHULUAN
%================================================

\begin{frame}
    \frametitle{Capaian Pembelajaran}
    \framesubtitle{Yang bisa kamu lakukan setelah sesi ini}
    \small
    \begin{enumerate}
        \item Memahami mengapa \textit{Deep Learning} lebih unggul dari ML tradisional untuk teks.
        \item Mengenal arsitektur \textbf{RNN, LSTM, dan BiLSTM} beserta konsep \textit{attention}.
        \item Memahami cara kerja \textbf{Transformer} dan \textit{fine-tuning} model pre-trained.
        \item Mampu melihat perbandingan performa antar model dari kode yang sudah berjalan.
    \end{enumerate}
    \begin{exampleblock}{Fokus Kelas Ini}
        Penjelasan lebih banyak \textbf{menggunakan intuisi visual} daripada rumus matematika yang rumit.
    \end{exampleblock}
\end{frame}

\begin{frame}
    \frametitle{Mengapa Deep Learning?}
    \framesubtitle{Keterbatasan yang membuat kita perlu naik level}
    \small
    Pada sesi-sesi sebelumnya kita sudah pakai \textbf{TF-IDF + SVM / PyCaret}.\\[0.4em]
    Ternyata ada beberapa hal yang \tem{tidak bisa} dilakukan pendekatan itu:
    \vspace{0.3em}
    \begin{itemize}
        \item \textbf{Tidak peduli urutan kata} — "saya \textit{tidak} suka" vs "saya suka \textit{tidak} ada masalah" dianggap mirip.
        \item \textbf{Tidak paham konteks} — kata "bank" di kalimat "tepi sungai" vs "menabung di bank" dianggap sama.
        \item \textbf{Dimensi fitur meledak} — setiap kata unik menjadi satu kolom; ribuan kata = ribuan kolom.
    \end{itemize}
\end{frame}

\begin{frame}
    \frametitle{Mengapa Deep Learning?}
    \framesubtitle{Apa yang ditawarkan Deep Learning}
    \small
    \textit{Deep Learning} menyelesaikan ketiga masalah tadi dengan cara:
    \vspace{0.3em}
    \begin{itemize}
        \item \textbf{Belajar representasi sendiri} — model secara otomatis menemukan pola penting dari data, tanpa kita perlu rekayasa fitur manual.
        \item \textbf{Memori urutan} — arsitektur RNN / LSTM "membaca" kata satu per satu sambil mengingat apa yang sudah dibaca.
        \item \textbf{Konteks global} — mekanisme \textit{attention} memungkinkan model melihat seluruh kalimat sekaligus.
    \end{itemize}
    \begin{exampleblock}{Analogi Sederhana}
        Kalau TF-IDF seperti \textbf{kamus kata} yang hanya menghitung frekuensi, maka DL seperti \textbf{orang yang membaca} — ia mengerti konteks dan urutan.
    \end{exampleblock}
\end{frame}

\begin{frame}
    \frametitle{Disclaimer Dataset}
    \begin{alertblock}{Peringatan Konten Sensitif}
        Dataset yang digunakan pada sesi ini berisi teks seputar \textbf{kesehatan mental} — termasuk topik depresi, kecemasan, dan pikiran bunuh diri.
        \vspace{0.25em}

        Dataset ini diambil dari sumber publik dan digunakan \textbf{semata-mata untuk tujuan edukasi} penelitian Pemrosesan Bahasa Alami.
        \vspace{0.25em}

        Jika kamu merasa terganggu dengan konten ini, silakan sampaikan ke pengajar.
    \end{alertblock}
\end{frame}

%================================================
% SECTION 1: REKAP & JEMBATAN
%================================================
\section{Rekap \& Jembatan dari Sesi Sebelumnya}

\begin{frame}
    \frametitle{Pipeline Traditional ML}
    \framesubtitle{Cara kerja yang sudah kita kenal}
    \small
    Setiap langkah dikerjakan \textbf{secara manual} oleh programmer:

    \vspace{0.6em}
    \begin{center}
    \begin{tikzpicture}[
        node distance=0.55cm and 0.35cm,
        box/.style={
            draw=nyubluedarker, fill=nyublue!12,
            rounded corners=3pt,
            minimum height=0.75cm, text width=1.85cm,
            align=center, font=\footnotesize
        },
        arr/.style={-{Stealth[length=5pt]}, thick, color=nyubluedarker}
    ]
        \node[box] (raw)   {Teks Mentah};
        \node[box, right=of raw]  (prep)  {Preprocessing \\ \tiny(lowercase, stopword)};
        \node[box, right=of prep] (tfidf) {TF-IDF \\ \tiny(matriks fitur)};
        \node[box, right=of tfidf](model) {Classifier \\ \tiny(SVM / RF)};
        \node[box, right=of model](pred)  {Prediksi \\ \tiny(label)};

        \draw[arr] (raw)   -- (prep);
        \draw[arr] (prep)  -- (tfidf);
        \draw[arr] (tfidf) -- (model);
        \draw[arr] (model) -- (pred);
    \end{tikzpicture}
    \end{center}

    \vspace{0.5em}
    \begin{itemize}
        \item \textbf{Fitur dibuat manual}: kita yang memilih cara merepresentasikan teks (TF-IDF, BoW).
        \item \textbf{Model hanya belajar dari fitur tersebut} — bukan dari teks aslinya.
    \end{itemize}
\end{frame}

\begin{frame}
    \frametitle{Pipeline Deep Learning}
    \framesubtitle{Cara kerja yang akan kita pelajari hari ini}
    \small
    Model \textbf{belajar sendiri} cara terbaik merepresentasikan teks:

    \vspace{0.6em}
    \begin{center}
    \begin{tikzpicture}[
        node distance=0.5cm and 0.3cm,
        box/.style={
            draw=nyupurple, fill=nyupurple!10,
            rounded corners=3pt,
            minimum height=0.75cm, text width=1.85cm,
            align=center, font=\footnotesize
        },
        arr/.style={-{Stealth[length=5pt]}, thick, color=nyupurple}
    ]
        \node[box] (raw)   {Teks Mentah};
        \node[box, right=of raw]  (tok)   {Tokenisasi \\ \tiny(angka indeks)};
        \node[box, right=of tok]  (emb)   {Embedding \\ \tiny(vektor makna)};
        \node[box, right=of emb]  (nn)    {Neural Net \\ \tiny(RNN/LSTM/BERT)};
        \node[box, right=of nn]   (pred)  {Prediksi \\ \tiny(label)};

        \draw[arr] (raw)  -- (tok);
        \draw[arr] (tok)  -- (emb);
        \draw[arr] (emb)  -- (nn);
        \draw[arr] (nn)   -- (pred);
    \end{tikzpicture}
    \end{center}

    \vspace{0.5em}
    \begin{itemize}
        \item \textbf{Embedding} menggantikan TF-IDF — representasi kata \textit{dipelajari} oleh model.
        \item \textbf{Neural Network} sekaligus menjadi ekstraktor fitur \textit{dan} klasifier.
    \end{itemize}
\end{frame}

\begin{frame}
    \frametitle{Perbandingan Pipeline: ML vs.\ DL}
    \framesubtitle{Di mana letak perbedaan utamanya?}
    \small
    \begin{columns}[T]
        \column{0.48\textwidth}
        \begin{block}{Traditional ML}
            \begin{itemize}
                \item Fitur \textbf{dibuat manual} (TF-IDF)
                \item Tidak peduli \textbf{urutan kata}
                \item Interpretasi lebih \textbf{mudah}
                \item Butuh \textbf{data lebih sedikit}
                \item Training \textbf{cepat}
            \end{itemize}
        \end{block}

        \column{0.48\textwidth}
        \begin{exampleblock}{Deep Learning}
            \begin{itemize}
                \item Fitur \textbf{dipelajari otomatis}
                \item Memahami \textbf{urutan \& konteks}
                \item Model lebih \textbf{kompleks}
                \item Butuh \textbf{data lebih banyak}
                \item Training \textbf{lebih lama}
            \end{itemize}
        \end{exampleblock}
    \end{columns}

    \vspace{0.6em}
    \begin{center}
        \small
        \textcolor{nyubluedarker}{%
            $\Rightarrow$ Tidak ada yang selalu lebih baik — pilihan tergantung data, waktu, dan kebutuhan.%
        }
    \end{center}
\end{frame}

\begin{frame}
    \frametitle{Dari Fitur Manual ke Representasi Terpelajar}
    \framesubtitle{Lompatan konseptual terbesar dalam Deep Learning}
    \small
    \begin{columns}[T]
        \column{0.46\textwidth}
        \textbf{TF-IDF (Fitur Manual)}\\[0.3em]
        Setiap kata menjadi satu kolom angka.\\[0.3em]
        \begin{center}
        \begin{tabular}{lcc}
            \toprule
            Kata & dok.1 & dok.2 \\
            \midrule
            \texttt{sedih}  & 0.42 & 0.00 \\
            \texttt{cemas}  & 0.00 & 0.55 \\
            \texttt{senang} & 0.10 & 0.33 \\
            \bottomrule
        \end{tabular}
        \end{center}
        \footnotesize Kata yang mirip maknanya \textbf{tidak} saling terkait.

        \column{0.50\textwidth}
        \textbf{Word Embedding (Terpelajar)}\\[0.3em]
        Setiap kata menjadi \textbf{vektor kontinu} berdimensi kecil.\\[0.3em]
        \begin{center}
        \begin{tabular}{lcc}
            \toprule
            Kata & dim-1 & dim-2 \\
            \midrule
            \texttt{sedih}  & $-$0.8 & $+$0.3 \\
            \texttt{cemas}  & $-$0.7 & $+$0.4 \\
            \texttt{senang} & $+$0.9 & $-$0.2 \\
            \bottomrule
        \end{tabular}
        \end{center}
        \footnotesize \texttt{sedih} dan \texttt{cemas} \textbf{berdekatan} di ruang vektor karena maknanya mirip.
    \end{columns}
\end{frame}

\begin{frame}
    \frametitle{Dataset: Dari Sesi 1 ke Sesi Ini}
    \framesubtitle{Mengapa kita ganti dataset?}
    \small
    \begin{columns}[T]
        \column{0.46\textwidth}
        \begin{block}{Sesi 1 — Indonesian Chat}
            \begin{itemize}
                \item Bahasa Indonesia
                \item 2 kelas: \textit{toxic} vs \textit{non-toxic}
                \item Klasifikasi sederhana
                \item Cocok untuk ML tradisional
            \end{itemize}
        \end{block}

        \column{0.50\textwidth}
        \begin{exampleblock}{Sesi Ini — Mental Health Sentiment}
            \begin{itemize}
                \item Bahasa Inggris
                \item \textbf{7 kelas} emosi/kondisi mental
                \item Nuansa makna sangat halus
                \item Membutuhkan pemahaman \textbf{konteks}
            \end{itemize}
        \end{exampleblock}
    \end{columns}

    \vspace{0.6em}
    Dataset dengan 7 kelas yang mirip maknanya menjadi \textbf{alasan kuat} mengapa kita perlu metode yang lebih canggih dari TF-IDF.
\end{frame}

\begin{frame}
    \frametitle{Sekilas Dataset: Mental Health Sentiment}
    \framesubtitle{7 Kelas Label yang Akan Kita Prediksi}
    \footnotesize
    {\renewcommand{\arraystretch}{1.15}
    \begin{center}
    \begin{tabular}{clp{5.8cm}}
        \toprule
        \textbf{\#} & \textbf{Label} & \textbf{Deskripsi Singkat} \\
        \midrule
        1 & \texttt{Normal}              & Teks tanpa indikasi kondisi mental tertentu \\
        2 & \texttt{Depression}          & Ekspresi rasa sedih, putus asa, tidak berdaya \\
        3 & \texttt{Suicidal}            & Pikiran atau niatan menyakiti diri \\
        4 & \texttt{Anxiety}             & Kekhawatiran berlebih, rasa takut, panik \\
        5 & \texttt{Bipolar}             & Perubahan suasana hati ekstrem (senang$\leftrightarrow$sedih) \\
        6 & \texttt{Stress}              & Tekanan hidup, beban pikiran \\
        7 & \texttt{Personality Disorder}& Pola pikir atau perilaku yang tidak biasa \\
        \bottomrule
    \end{tabular}
    \end{center}}
    \textcolor{nyuorange}{\textbf{Tantangan:}} Banyak kelas yang sulit dibedakan — misalnya \texttt{Depression} vs \texttt{Suicidal}. Inilah yang membuat tugasnya menarik!
\end{frame}

%================================================
% SECTION 2: DASAR-DASAR DEEP LEARNING
%================================================
\section{Dasar-Dasar Deep Learning}

% --- 2.1 Neuron Buatan ---
\begin{frame}
    \frametitle{Neuron: Unit Terkecil Deep Learning}
    \framesubtitle{Meniru cara kerja sel saraf di otak}
    \small
    Bayangkan neuron seperti \textbf{saklar cerdas}: ia menerima beberapa sinyal,
    menimbang-nimbangnya, lalu memutuskan apakah harus "menyala" atau tidak.

    \vspace{0.3em}
    \begin{center}
    \begin{tikzpicture}[
        node distance=0.4cm and 1.1cm,
        inp/.style={circle, draw=nyubluedarker, fill=nyublue!15, minimum size=0.7cm, font=\footnotesize},
        neu/.style={circle, draw=nyupurple, fill=nyupurple!15, minimum size=1cm, font=\footnotesize, align=center},
        outnd/.style={circle, draw=nyuorange!80!black, fill=nyuorange!15, minimum size=0.7cm, font=\footnotesize},
        arr/.style={-{Stealth[length=4.5pt]}, thick}
    ]
        \node[inp] (x1) {$x_1$};
        \node[inp, below=of x1] (x2) {$x_2$};
        \node[inp, below=of x2] (x3) {$x_3$};

        \node[neu, right=1.4cm of x2] (n) {$\Sigma$\\+\\$f$};

        \node[outnd, right=1.4cm of n] (y) {$y$};

        \draw[arr, color=nyubluedarker] (x1) -- node[above, font=\tiny, sloped]{$w_1$} (n);
        \draw[arr, color=nyubluedarker] (x2) -- node[above, font=\tiny]{$w_2$} (n);
        \draw[arr, color=nyubluedarker] (x3) -- node[above, font=\tiny, sloped]{$w_3$} (n);
        \draw[arr, color=nyupurple] (n) -- (y);

        % bias annotation
        \node[above=0.15cm of n, font=\tiny, color=nyupurple] {$+\,b$};
    \end{tikzpicture}
    \end{center}

    \vspace{0.2em}
    \footnotesize
    \begin{itemize}
        \item $x_1, x_2, x_3$ = \textbf{input}, $\;$ $w_1, w_2, w_3$ = \textbf{weight} (bobot kepentingan)
        \item $b$ = \textbf{bias} (pergeseran ambang), $\;$ $f$ = \textbf{fungsi aktivasi}
        \item Rumus: $y = f\bigl(w_1 x_1 + w_2 x_2 + w_3 x_3 + b\bigr)$
    \end{itemize}
\end{frame}

\begin{frame}
    \frametitle{Weight \& Bias: Parameter yang Dipelajari}
    \framesubtitle{Model belajar = model menyetel angka-angka ini}
    \footnotesize
    \begin{columns}[T]
        \column{0.48\textwidth}
        \begin{block}{Weight ($w$)}
            \begin{itemize}
                \item Menentukan \textbf{seberapa besar pengaruh} suatu input
                \item $w$ besar $\rightarrow$ sangat berpengaruh
                \item $w$ negatif $\rightarrow$ bersifat "menghambat"
            \end{itemize}
        \end{block}

        \column{0.48\textwidth}
        \begin{exampleblock}{Bias ($b$)}
            \begin{itemize}
                \item Menggeser titik keputusan neuron
                \item Neuron bisa aktif meski semua input = 0
                \item Seperti $b$ dalam $y = mx + b$
            \end{itemize}
        \end{exampleblock}
    \end{columns}

    \vspace{0.3em}
    \begin{alertblock}{Inisialisasi Weight}
        Dimulai dari nilai acak kecil. Metode umum: \textbf{Xavier} atau \textbf{He initialization}. Nilai disesuaikan bertahap selama training.
    \end{alertblock}
\end{frame}

% --- 2.2 Fungsi Aktivasi ---
\begin{frame}
    \frametitle{Fungsi Aktivasi}
    \framesubtitle{Memberi kemampuan "berpikir non-linear" pada neuron}
    \small
    Tanpa fungsi aktivasi, seluruh jaringan hanya melakukan operasi linear.

    \vspace{0.2em}
    \begin{center}
    {\footnotesize
    \begin{tabular}{llp{5.2cm}}
        \toprule
        \textbf{Fungsi} & \textbf{Rentang Output} & \textbf{Kapan digunakan} \\
        \midrule
        \textbf{ReLU}    & $[0, \infty)$  & Hidden layer — paling umum, komputasi cepat \\
        \textbf{Sigmoid} & $(0, 1)$       & Output layer klasifikasi \textit{biner} \\
        \textbf{Tanh}    & $(-1, 1)$      & Hidden layer RNN / LSTM — mengatasi batasan Sigmoid \\
        \textbf{Softmax} & $(0, 1)$, jumlah = 1 & Output layer klasifikasi \textit{multikelas} \\
        \bottomrule
    \end{tabular}}
    \end{center}

    \vspace{0.15em}
    \begin{exampleblock}{Analogi ReLU}
        ReLU ($\max(0, x)$) seperti lampu yang hanya menyala saat ada listrik positif — nilai negatif $\rightarrow$ tetap 0.
    \end{exampleblock}
\end{frame}

\begin{frame}
    \frametitle{Grafik Fungsi Aktivasi}
    \framesubtitle{Visualisasi singkat bentuk masing-masing fungsi}
    \small
    \begin{center}
    \begin{tikzpicture}[scale=0.78]
        % ReLU
        \begin{scope}[xshift=0cm]
            \draw[->] (-1.2,0)--(1.5,0) node[right,font=\tiny]{$x$};
            \draw[->] (0,-0.3)--(0,1.5) node[above,font=\tiny]{$y$};
            \draw[thick, color=nyubluedarker] (-1.2,0)--(0,0)--(1.2,1.2);
            \node[font=\scriptsize, below=2pt] at (0.15,-0.3) {\textbf{ReLU}};
        \end{scope}
        % Sigmoid
        \begin{scope}[xshift=3.5cm]
            \draw[->] (-1.5,0)--(1.5,0) node[right,font=\tiny]{$x$};
            \draw[->] (0,-0.1)--(0,1.3) node[above,font=\tiny]{$y$};
            \draw[thick, color=nyupurple, domain=-1.4:1.4, samples=40]
                plot (\x, {1/(1+exp(-3*\x))});
            \node[font=\scriptsize, below=2pt] at (0.15,-0.3) {\textbf{Sigmoid}};
        \end{scope}
        % Tanh
        \begin{scope}[xshift=7cm]
            \draw[->] (-1.5,0)--(1.5,0) node[right,font=\tiny]{$x$};
            \draw[->] (0,-1.3)--(0,1.3) node[above,font=\tiny]{$y$};
            \draw[thick, color=nyuorange!80!black, domain=-1.4:1.4, samples=40]
                plot (\x, {tanh(2.5*\x)});
            \node[font=\scriptsize, below=2pt] at (0.15,-1.5) {\textbf{Tanh}};
        \end{scope}
    \end{tikzpicture}
    \end{center}

    \vspace{0.2em}
    \footnotesize
    \textbf{Softmax} tidak digambar di sini karena ia bekerja pada \textit{vektor} (bukan satu angka) — mengubah sekumpulan skor menjadi distribusi probabilitas yang totalnya = 1.
\end{frame}

% --- 2.3 Layer ---
\begin{frame}
    \frametitle{Layer: Menyusun Neuron Menjadi Jaringan}
    \framesubtitle{Konsep \textit{depth} dalam Deep Learning}
    \small
    \begin{center}
    \begin{tikzpicture}[scale=0.72,
        node distance=0.3cm and 1.6cm,
        nd/.style={circle, draw, minimum size=0.45cm, font=\tiny},
        inp/.style={nd, draw=nyubluedarker, fill=nyublue!15},
        hid/.style={nd, draw=nyupurple, fill=nyupurple!12},
        outnd/.style={nd, draw=nyuorange!80!black, fill=nyuorange!15},
        arr/.style={-{Stealth[length=3.5pt]}, gray, thin}
    ]
        % Input layer
        \foreach \i in {1,...,3} {
            \node[inp] (i\i) at (0, -\i*0.65) {};
        }
        % Hidden layer 1
        \foreach \i in {1,...,4} {
            \node[hid] (h1\i) at (2, -\i*0.55+0.1) {};
        }
        % Hidden layer 2
        \foreach \i in {1,...,4} {
            \node[hid] (h2\i) at (4, -\i*0.55+0.1) {};
        }
        % Output layer
        \foreach \i in {1,...,2} {
            \node[outnd] (o\i) at (6, -\i*0.65) {};
        }
        % Connections
        \foreach \i in {1,...,3} {
            \foreach \j in {1,...,4} { \draw[arr] (i\i) -- (h1\j); }
        }
        \foreach \i in {1,...,4} {
            \foreach \j in {1,...,4} { \draw[arr] (h1\i) -- (h2\j); }
        }
        \foreach \i in {1,...,4} {
            \foreach \j in {1,...,2} { \draw[arr] (h2\i) -- (o\j); }
        }
        % Labels
        \node[font=\tiny, color=nyubluedarker, below=2pt] at (0,-2.2) {\textbf{Input}};
        \node[font=\tiny, color=nyupurple, below=2pt]     at (2,-2.3) {\textbf{Hidden 1}};
        \node[font=\tiny, color=nyupurple, below=2pt]     at (4,-2.3) {\textbf{Hidden 2}};
        \node[font=\tiny, color=nyuorange!80!black, below=2pt] at (6,-1.5) {\textbf{Output}};
    \end{tikzpicture}
    \end{center}

    \footnotesize
    \begin{itemize}
        \item \textbf{Input layer}: menerima data mentah (teks yang sudah di-encode)
        \item \textbf{Hidden layer}: mengekstrak pola — makin banyak layer, makin \textit{dalam}
        \item \textbf{Output layer}: menghasilkan prediksi (7 kelas $\rightarrow$ 7 neuron output)
    \end{itemize}
\end{frame}

% --- 2.4 Loss Function ---
\begin{frame}
    \frametitle{Loss Function: Mengukur Seberapa Salah Model}
    \framesubtitle{Sinyal yang memberi tahu model "kamu masih keliru"}
    \small
    \textbf{Loss function} mengukur selisih prediksi model vs.~jawaban benar.
    Makin kecil \textit{loss}, makin bagus.

    \begin{exampleblock}{Cross-Entropy Loss — Klasifikasi Multikelas}
        \[
            \mathcal{L} = -\sum_{c=1}^{C} y_c \cdot \log(\hat{p}_c)
        \]
        \footnotesize
        \begin{itemize}
            \item $y_c = 1$ jika kelas $c$ benar, 0 jika tidak
            \item $\hat{p}_c$ = probabilitas prediksi untuk kelas $c$ (output Softmax)
        \end{itemize}
    \end{exampleblock}

    \vspace{0.15em}
    \begin{alertblock}{Contoh}
        \footnotesize Jawaban: \texttt{Anxiety}. Prediksi \texttt{Anxiety} 90\% $\rightarrow$ loss kecil. Prediksi \texttt{Normal} 90\% $\rightarrow$ loss besar.
    \end{alertblock}
\end{frame}

% --- 2.5 Feedforward / MLP ---
\begin{frame}
    \frametitle{Feedforward Network (MLP)}
    \framesubtitle{Neural network paling sederhana — data mengalir satu arah}
    \small
    \textbf{MLP} (\textit{Multi-Layer Perceptron}) adalah jaringan saraf di mana data bergerak \textbf{hanya dari kiri ke kanan} — tidak ada putaran balik.

    \vspace{0.4em}
    \begin{center}
    \begin{tikzpicture}[
        box/.style={draw=nyubluedarker, fill=nyublue!10, rounded corners=3pt,
                    minimum height=0.8cm, text width=2cm, align=center, font=\footnotesize},
        arr/.style={-{Stealth[length=5pt]}, thick, color=nyubluedarker}
    ]
        \node[box] (inp)  {Input\\(teks angka)};
        \node[box, right=0.8cm of inp]  (h1) {Hidden\\Layer 1};
        \node[box, right=0.8cm of h1]   (h2) {Hidden\\Layer 2};
        \node[box, right=0.8cm of h2]   (out){Output\\(7 kelas)};

        \draw[arr] (inp) -- (h1);
        \draw[arr] (h1)  -- (h2);
        \draw[arr] (h2)  -- (out);

        % annotations
        \node[font=\tiny, above=3pt of h1, color=nyupurple]{ReLU};
        \node[font=\tiny, above=3pt of h2, color=nyupurple]{ReLU};
        \node[font=\tiny, above=3pt of out, color=nyuorange!80!black]{Softmax};
    \end{tikzpicture}
    \end{center}

    \vspace{0.3em}
    \begin{columns}[T]
        \column{0.48\textwidth}
        \textbf{Kelebihan MLP:}
        \begin{itemize}
            \item Sederhana dan cepat dilatih
            \item Cocok untuk data tabular
        \end{itemize}

        \column{0.48\textwidth}
        \textbf{Kekurangan untuk Teks:}
        \begin{itemize}
            \item Tidak mempertimbangkan \textbf{urutan kata}
            \item Inilah yang nanti diselesaikan RNN
        \end{itemize}
    \end{columns}
\end{frame}

% --- 2.6 Overfitting & Regularisasi ---
\begin{frame}
    \frametitle{Overfitting: Model Terlalu "Hafal"}
    \framesubtitle{Masalah klasik yang sering muncul saat melatih model}
    \small
    \textbf{Overfitting} terjadi ketika model terlalu fokus menghafal data latih,
    sehingga performa di data baru (data uji) justru buruk.

    \vspace{0.4em}
    \begin{center}
    \begin{tikzpicture}[scale=0.78]
        % axes
        \draw[->] (0,0) -- (5.5,0) node[right,font=\footnotesize]{Epoch};
        \draw[->] (0,0) -- (0,2.8) node[above,font=\footnotesize]{Loss};

        % train loss (goes down smoothly)
        \draw[thick, color=nyubluedarker, domain=0.3:5, samples=50]
            plot (\x, {2.2*exp(-0.7*\x)+0.15});
        \node[font=\tiny, color=nyubluedarker] at (4.5,0.4) {Train Loss};

        % val loss (goes down then up)
        \draw[thick, color=nyuorange!80!black, domain=0.3:5, samples=50]
            plot (\x, {2.2*exp(-0.7*\x)+0.15 + 0.5*max(0,\x-2.5)^1.2});
        \node[font=\tiny, color=nyuorange!80!black] at (4.2,2.0) {Val Loss};

        % overfitting marker
        \draw[dashed, gray] (2.5,0) -- (2.5,2.5);
        \node[font=\tiny, gray, above] at (2.5,2.5) {mulai \textit{overfit}};
    \end{tikzpicture}
    \end{center}

    \vspace{0.1em}
    \footnotesize Saat \textit{val loss} mulai naik sementara \textit{train loss} terus turun, itulah tanda overfitting.
\end{frame}

\begin{frame}
    \frametitle{Regularisasi: Cara Mencegah Overfitting}
    \framesubtitle{Tiga teknik yang paling sering digunakan}
    \small
    \begin{columns}[T]
        \column{0.32\textwidth}
        \begin{block}{Dropout}
            Saat training, sebagian neuron \textbf{dimatikan secara acak} di setiap iterasi.\\[0.3em]
            Memaksa model tidak bergantung pada neuron tertentu saja.\\[0.3em]
            \footnotesize Contoh: \texttt{Dropout(0.3)} $\rightarrow$ 30\% neuron dimatikan tiap langkah.
        \end{block}

        \column{0.32\textwidth}
        \begin{exampleblock}{L2 Regularization}
            Menambahkan \textbf{penalti} pada nilai weight yang terlalu besar ke dalam loss function.\\[0.3em]
            Weight besar $\rightarrow$ model mengandalkan fitur tertentu secara berlebihan.\\[0.3em]
            \footnotesize Disebut juga \textit{weight decay}.
        \end{exampleblock}

        \column{0.32\textwidth}
        \begin{alertblock}{Early Stopping}
            Hentikan training \textbf{saat val loss mulai naik}, bukan saat mencapai epoch maksimum.\\[0.3em]
            \footnotesize Parameter: \texttt{patience=3} $\rightarrow$ training berhenti jika val loss tidak membaik selama 3 epoch berturut-turut.
        \end{alertblock}
    \end{columns}
\end{frame}

\end{document}
